\begin{animation}[id="segtree-query", label="Segment Tree Range Sum Query [2,5]"]
\shape{st}{Tree}{data=[3,1,4,1,5,9,2,6], kind="segtree", show_sum=true}

\step
\narrate{Segment tree built from array [3,1,4,1,5,9,2,6]. We want the sum of range [2,5].}

\step
\recolor{st.node["[0,7]"]}{state=current}
\narrate{Start at root [0,7] with sum 31. Range [2,5] partially overlaps, so recurse both children.}

\step
\recolor{st.node["[0,7]"]}{state=dim}
\recolor{st.node["[0,3]"]}{state=current}
\recolor{st.node["[4,7]"]}{state=current}
\narrate{Visit [0,3] (sum=9) and [4,7] (sum=22). Both partially overlap [2,5].}

\step
\recolor{st.node["[0,3]"]}{state=dim}
\recolor{st.node["[4,7]"]}{state=dim}
\recolor{st.node["[0,1]"]}{state=done}
\recolor{st.node["[2,3]"]}{state=good}
\recolor{st.node["[4,5]"]}{state=good}
\recolor{st.node["[6,7]"]}{state=done}
\narrate{[0,1] is outside range, skip. [2,3] (sum=5) fully inside, take it. [4,5] (sum=14) fully inside, take it. [6,7] is outside, skip.}

\step
\recolor{st.node["[2,3]"]}{state=current}
\recolor{st.node["[4,5]"]}{state=current}
\narrate{Answer = 5 + 14 = 19. The sum of elements at indices 2 through 5 is 19.}
\end{animation}
